\chapter{2024年全国甲卷（理）}
\sloppy
\hfuzz=120pt

\section{单选题}

\begin{problem}
设 \( z=5+\i \)，则 \( \i \left( \bar z+z \right)= \)（\quad）

\choices
    {\( 10\i \)}
    {\( 2\i \)}
    {\( 10 \)}
    {\( -2 \)}

\end{problem}

\begin{answer}
A
\end{answer}

\begin{solution}
根据 \( z=5+\i \)，可得 \( \bar z=5-\i \)，所以
\(z+\bar z=10\)，
则
\(\i \left( \bar z+z \right)=10\i\).

故选 A.
\end{solution}

\begin{problem}
集合 \( A=\{1,2,3,4,5,9\} \)，\( B=\left\{ x \mid \sqrt{x}\in A \right\} \)，则 \( C_A\left( A\cap B \right)= \)（\quad）

\choices
    {\( \{1,4,9\} \)}
    {\( \{3,4,9\} \)}
    {\( \{1,2,3\} \)}
    {\( \{2,3,5\} \)}

\end{problem}

\begin{answer}
D
\end{answer}

\begin{solution}
因为 \( A=\{1,2,3,4,5,9\} \)，\( B=\left\{ x \mid \sqrt{x}\in A \right\} \)，所以
\(B=\{1,4,9,16,25,81\}\)，
则
\(A\cap B=\{1,4,9\}\)，
故
\(C_A\left( A\cap B \right)=\{2,3,5\}\).

故选 D.
\end{solution}

\begin{problem}
若实数 \( x,y \) 满足约束条件 \( 4x-3y-3\ge 0 \)，\( x-2y-2\le 0 \)，\( 2x+6y-9\le 0 \)，则 \( z=x-5y \) 的最小值为（\quad）

\begin{center}
\bitmapinclude[max width=6cm,max totalheight=4cm]{img/2024/national_paper_A_science/q3_fig1.png}
\end{center}

\choices
    {\( 5 \)}
    {\( \dfrac{1}{2} \)}
    {\( -2 \)}
    {\( -\dfrac{7}{2} \)}

\end{problem}

\begin{answer}
D
\end{answer}

\begin{solution}
实数 \( x,y \) 满足 \( 4x-3y-3\ge 0 \)，\( x-2y-2\le 0 \)，\( 2x+6y-9\le 0 \)，作出可行域如图.

根据 \( z=x-5y \) 可得
\(y=\frac{1}{5}x-\frac{1}{5}z\)，
即 \( z \) 的几何意义为 \( y=\frac{1}{5}x-\frac{1}{5}z \) 的截距的 \( -5 \) 倍.

则该直线截距取最大值时，\( z \) 有最小值，此时直线 \( y=\frac{1}{5}x-\frac{1}{5}z \) 过点 \( A \).

联立 \( 4x-3y-3=0 \) 与 \( 2x+6y-9=0 \)，解得 \( x=\dfrac{3}{2} \)，\( y=1 \)，即 \( A\left( \dfrac{3}{2},1 \right) \)，则
\(z_{\min}=\frac{3}{2}-5\times 1=-\frac{7}{2}\).

故选 D.
\end{solution}

\begin{problem}
等差数列 \( \{a_n\} \) 的前 \( n \) 项和为 \( S_n \)，若 \( S_5=S_{10} \)，\( a_5=1 \)，则 \( a_1= \)（\quad）

\choices
    {\( -2 \)}
    {\( \dfrac{7}{3} \)}
    {\( 1 \)}
    {\( 2 \)}

\end{problem}

\begin{answer}
B
\end{answer}

\begin{solution}
由
\[
S_{10}-S_5=a_6+a_7+a_8+a_9+a_{10}=5a_8=0
\]
可得 \( a_8=0 \).

则等差数列 \( \{a_n\} \) 的公差 \( d=\frac{a_8-a_5}{3}=-\frac{1}{3} \)，故
\[
a_1=a_5-4d=1-4\times \left( -\frac{1}{3} \right)=\frac{7}{3}
\]

故选 B.
\end{solution}

\begin{problem}
已知双曲线 \( C:\dfrac{y^2}{a^2}-\dfrac{x^2}{b^2}=1 \)（\( a>0 \)，\( b>0 \)）的上、下焦点分别为 \( F_1(0,4) \)，\( F_2(0,-4) \)，点 \( P(-6,4) \) 在该双曲线上，则该双曲线的离心率为（\quad）

\choices
    {\( 4 \)}
    {\( 3 \)}
    {\( 2 \)}
    {\( \sqrt{2} \)}

\end{problem}

\begin{answer}
C
\end{answer}

\begin{solution}
根据题意，\( F_1(0,-4) \)，\( F_2(0,4) \)，\( P(-6,4) \)，则
\(\left| F_1F_2 \right|=2c=8\)，
\(\left| PF_1 \right|=\sqrt{6^2+\left( 4+4 \right)^2}=10\)，
\(\left| PF_2 \right|=\sqrt{6^2+\left( 4-4 \right)^2}=6\).

由双曲线定义可得
\(2a=\left| PF_1 \right|-\left| PF_2 \right|=10-6=4\)，
则
\(e=\frac{2c}{2a}=\frac{8}{4}=2\).

故选 C.
\end{solution}

\begin{problem}
设函数 \( f(x)=\dfrac{\e^x+2\sin x}{1+x^2} \)，则曲线 \( y=f(x) \) 在 \( (0,1) \) 处的切线与两坐标轴围成的三角形的面积为（\quad）

\choices
    {\( \dfrac{1}{6} \)}
    {\( \dfrac{1}{3} \)}
    {\( \dfrac{1}{2} \)}
    {\( \dfrac{2}{3} \)}

\end{problem}

\begin{answer}
A
\end{answer}

\begin{solution}
借助导数的几何意义计算可得其在点 \( (0,1) \) 处的切线方程.

\[
f'(x)=\frac{\left( \e^x+2\cos x \right)\left( 1+x^2 \right)-\left( \e^x+2\sin x \right)\cdot 2x}{\left( 1+x^2 \right)^2}
\]

所以
\[
f'(0)=\frac{\left( \e^0+2\cos 0 \right)\left( 1+0 \right)-\left( \e^0+2\sin 0 \right)\times 0}{\left( 1+0 \right)^2}=3
\]

即该切线方程为
\[
y-1=3x
\]
也就是
\[
y=3x+1
\]

令 \( x=0 \)，则 \( y=1 \)；令 \( y=0 \)，则 \( x=-\dfrac{1}{3} \).

所以该切线与两坐标轴所围成的三角形面积
\(S=\frac{1}{2}\times 1\times \left| -\frac{1}{3} \right|=\frac{1}{6}\).

故选 A.
\end{solution}

\begin{problem}
函数 \( f(x)=-x^2+\left( \e^x-\e^{-x} \right)\sin x \) 在区间 \( [-2.8,2.8] \) 的大致图像为（\quad）

\choices
    {\bitmapinclude[max width=4.0cm,max totalheight=3.1cm]{img/2024/national_paper_A_science/q7_fig1.png}}
    {\bitmapinclude[max width=4.0cm,max totalheight=3.1cm]{img/2024/national_paper_A_science/q7_fig2.png}}
    {\bitmapinclude[max width=4.0cm,max totalheight=3.1cm]{img/2024/national_paper_A_science/q7_fig3.png}}
    {\bitmapinclude[max width=4.0cm,max totalheight=3.1cm]{img/2024/national_paper_A_science/q7_fig4.png}}

\end{problem}

\begin{answer}
B
\end{answer}

\begin{solution}
由
\[
f(-x)=-x^2+\left( \e^{-x}-\e^x \right)\sin(-x)=-x^2+\left( \e^x-\e^{-x} \right)\sin x=f(x)
\]
可知函数定义域为 \( [-2.8,2.8] \) 时，\( f(x) \) 为偶函数，故 A、C 错误.

又
\[
f(1)=-1+\left( \e-\frac{1}{\e} \right)\sin 1>-1+\left( \e-\frac{1}{\e} \right)\sin \frac{\pi}{6}=\frac{\e}{2}-1-\frac{1}{2\e}>\frac{1}{4}-\frac{1}{2\e}>0
\]
故 D 错误.

故选 B.
\end{solution}

\begin{problem}
已知 \( \dfrac{\cos \alpha}{\cos \alpha-\sin \alpha}=\sqrt{3} \)，则 \( \tan \left( \alpha+\dfrac{\pi}{4} \right)= \)（\quad）

\choices
    {\( 2\sqrt{3}+1 \)}
    {\( 2\sqrt{3}-1 \)}
    {\( \dfrac{\sqrt{3}}{2} \)}
    {\( 1-\sqrt{3} \)}

\end{problem}

\begin{answer}
B
\end{answer}

\begin{solution}
因为
\(\frac{\cos \alpha}{\cos \alpha-\sin \alpha}=\sqrt{3}\)，
所以
\(\frac{1}{1-\tan \alpha}=\sqrt{3}\)，
即
\(\tan \alpha=1-\frac{\sqrt{3}}{3}\).

所以
\(\tan \left( \alpha+\frac{\pi}{4} \right)=\frac{\tan \alpha+1}{1-\tan \alpha}=2\sqrt{3}-1\).

故选 B.
\end{solution}

\begin{problem}
已知向量 \( \bs a=(x+1,x) \)，\( \bs b=(x,2) \)，则（\quad）

\choices
    {“\( x=-3 \)”是“\( \bs a\perp \bs b \)”的必要条件}
    {“\( x=-3 \)”是“\( \bs a//\bs b \)”的必要条件}
    {“\( x=0 \)”是“\( \bs a\perp \bs b \)”的充分条件}
    {“\( x=-1+\sqrt{3} \)”是“\( \bs a//\bs b \)”的充分条件}

\end{problem}

\begin{answer}
C
\end{answer}

\begin{solution}
A，当 \( \bs a\perp \bs b \) 时，则
\(\bs a\cdot \bs b=0\)，
所以
\(x(x+1)+2x=0\)，
解得 \( x=0 \) 或 \( x=-3 \)，即必要性不成立，A 错误.

B，当 \( \bs a//\bs b \) 时，则
\(2(x+1)=x^2\)，
解得
\(x=1\pm \sqrt{3}\)，
即必要性不成立，B 错误.

C，当 \( x=0 \) 时，
\(\bs a=(1,0)\)，\(\bs b=(0,2)\)，
故
\(\bs a\cdot \bs b=0\)，
所以 \( \bs a\perp \bs b \)，即充分性成立，C 正确.

D，当 \( x=-1+\sqrt{3} \) 时，不满足
\(2(x+1)=x^2\)，
所以 \( \bs a//\bs b \) 不成立，即充分性不成立，D 错误.

故选 C.
\end{solution}

\begin{problem}
设 \( \alpha,\beta \) 是两个平面，\( m,n \) 是两条直线，且 \( \alpha\cap\beta=m \). 下列四个命题：

\circled{1} 若 \( m//n \)，则 \( n//\alpha \) 或 \( n//\beta \)；

\circled{2} 若 \( m\perp n \)，则 \( n\perp \alpha \)，\( n\perp \beta \)；

\circled{3} 若 \( n//\alpha \)，且 \( n//\beta \)，则 \( m//n \)；

\circled{4} 若 \( n \) 与 \( \alpha \) 和 \( \beta \) 所成的角相等，则 \( m\perp n \).

其中所有真命题的编号是（\quad）

\choices
    {\( \circled{1}\circled{3} \)}
    {\( \circled{2}\circled{4} \)}
    {\( \circled{1}\circled{2}\circled{3} \)}
    {\( \circled{1}\circled{3}\circled{4} \)}

\end{problem}

\begin{answer}
A
\end{answer}

\begin{solution}
\circled{1}，当 \( n\subset \alpha \) 时，因为 \( m//n \)，\( m\subset \beta \)，则 \( n//\beta \)；

当 \( n\subset \beta \) 时，因为 \( m//n \)，\( m\subset \alpha \)，则 \( n//\alpha \)；

当 \( n \) 既不在 \( \alpha \) 也不在 \( \beta \) 内时，因为 \( m//n \)，\( m\subset \alpha \)，\( m\subset \beta \)，则 \( n//\alpha \) 且 \( n//\beta \)，故 \circled{1} 正确.

\circled{2}，若 \( m\perp n \)，则 \( n \) 与 \( \alpha,\beta \) 不一定垂直，故 \circled{2} 错误.

\circled{3}，过直线 \( n \) 分别作两平面与 \( \alpha,\beta \) 分别相交于直线 \( s \) 和直线 \( t \). 因为 \( n//\alpha \)，过直线 \( n \) 的平面与平面 \( \alpha \) 的交线为直线 \( s \)，则根据线面平行的性质定理知 \( n//s \). 又因为 \( n//\beta \)，过直线 \( n \) 的平面与平面 \( \beta \) 的交线为直线 \( t \)，则根据线面平行的性质定理知 \( n//t \)，则 \( s//t \). 因为 \( s\nsubseteq \beta \)，\( t\subset \beta \)，则 \( s//\beta \). 因为 \( s\subset \alpha \)，\( \alpha\cap\beta=m \)，则 \( s//m \)，又因为 \( n//s \)，则 \( m//n \)，故 \circled{3} 正确.

\circled{4}，若 \( \alpha\cap\beta=m \)，\( n \) 与 \( \alpha \) 和 \( \beta \) 所成的角相等，如果 \( n//\alpha \)，\( n//\beta \)，则 \( m//n \)，故 \circled{4} 错误.

\circled{1}\circled{3} 正确，故选 A.
\end{solution}

\begin{problem}
在 \( \triangle ABC \) 中内角 \( A,B,C \) 所对边分别为 \( a,b,c \)，若 \( B=\dfrac{\pi}{3} \)，\( b^2=\dfrac{9}{4}ac \)，则 \( \sin A+\sin C= \)（\quad）

\choices
    {\( \dfrac{3}{2} \)}
    {\( \sqrt{2} \)}
    {\( \dfrac{\sqrt{7}}{2} \)}
    {\( \dfrac{\sqrt{3}}{2} \)}

\end{problem}

\begin{answer}
C
\end{answer}

\begin{solution}
因为 \( B=\dfrac{\pi}{3} \)，\( b^2=\dfrac{9}{4}ac \)，由正弦定理得
\[
\sin A\sin C=\frac{4}{9}\sin^2 B=\frac{1}{3}
\]

由余弦定理可得
\(b^2=a^2+c^2-ac=\frac{9}{4}ac\)，
即
\(a^2+c^2=\frac{13}{4}ac\).

根据正弦定理得
\[
\sin^2 A+\sin^2 C=\frac{13}{4}\sin A\sin C=\frac{13}{12}
\]
所以
\[
\left( \sin A+\sin C \right)^2=\sin^2 A+\sin^2 C+2\sin A\sin C=\frac{7}{4}
\]

因为 \( A,C \) 为三角形内角，则 \( \sin A+\sin C>0 \)，故
\(\sin A+\sin C=\frac{\sqrt{7}}{2}\).

故选 C.
\end{solution}

\begin{problem}
已知 \( b \) 是 \( a,c \) 的等差中项，直线 \( ax+by+c=0 \) 与圆 \( x^2+y^2+4y-1=0 \) 交于 \( A,B \) 两点，则 \( \left| AB \right| \) 的最小值为（\quad）

\choices
    {\( 2 \)}
    {\( 3 \)}
    {\( 4 \)}
    {\( 2\sqrt{5} \)}

\end{problem}

\begin{answer}
C
\end{answer}

\begin{solution}
因为 \( a,b,c \) 成等差数列，所以 \( 2b=a+c \)，即 \( c=2b-a \).

代入直线方程 \( ax+by+c=0 \) 得
\(ax+by+2b-a=0\)，
即
\(a(x-1)+b(y+2)=0\).

令 \( x-1=0 \)，\( y+2=0 \)，得 \( x=1 \)，\( y=-2 \)，故直线恒过 \( (1,-2) \)，设 \( P(1,-2) \).

圆化为标准方程得
\(C:x^2+(y+2)^2=5\).

设圆心为 \( C \)，当 \( PC\perp AB \) 时，\( |AB| \) 最小.

因为
\(|PC|=1\)，\(|AC|=r=\sqrt{5}\)，
此时
\(|AB|=2|AP|=2\sqrt{AC^2-PC^2}=2\sqrt{5-1}=4\).

故选 C.
\end{solution}
\section{填空题}

\begin{problem}
\( \left( \dfrac{1}{3}+x \right)^{10} \) 的展开式中，各项系数的最大值是\fillinblank{}.

\end{problem}

\begin{answer}
5
\end{answer}

\begin{solution}
根据题意，展开式通项公式为
\(T_{r+1}=C_{10}^r\left( \frac{1}{3} \right)^{10-r}x^r\)，
其中 \( 0\le r\le 10 \)，且 \( r\in \Z \).

设展开式中第 \( r+1 \) 项系数最大，则 \( C_{10}^r\left( \dfrac{1}{3} \right)^{10-r}\ge C_{10}^{r+1}\left( \dfrac{1}{3} \right)^{9-r} \)，且 \( C_{10}^r\left( \dfrac{1}{3} \right)^{10-r}\ge C_{10}^{r-1}\left( \dfrac{1}{3} \right)^{11-r} \)，解得 \( \frac{29}{4}\le r\le \frac{33}{4} \).

又 \( r\in \Z \)，故 \( r=8 \)，所以展开式中系数最大的项是第 \( 9 \) 项，且该项系数为
\(C_{10}^8\left( \frac{1}{3} \right)^2=5\).

答案为：5.
\end{solution}

\begin{problem}
已知甲、乙两个圆台上、下底面的半径均为 \( r_1 \) 和 \( r_2 \)，母线长分别为 \( 2\left( r_2-r_1 \right) \) 和 \( 3\left( r_2-r_1 \right) \)，则两个圆台的体积之比 \( \dfrac{V_{\text{甲}}}{V_{\text{乙}}}= \)\fillinblank{}.

\end{problem}

\begin{answer}
\( \dfrac{\sqrt{6}}{4} \)
\end{answer}

\begin{solution}
根据题可得两个圆台的高分别为
\(h_{\text{甲}}=\sqrt{\left[ 2\left( r_1-r_2 \right) \right]^2-\left( r_1-r_2 \right)^2}=\sqrt{3}\left( r_1-r_2 \right)\)，
\(h_{\text{乙}}=\sqrt{\left[ 3\left( r_1-r_2 \right) \right]^2-\left( r_1-r_2 \right)^2}=2\sqrt{2}\left( r_1-r_2 \right)\).

所以
\[
\frac{V_{\text{甲}}}{V_{\text{乙}}}
=\frac{\dfrac{1}{3}\left( S_2+S_1+\sqrt{S_2S_1} \right)h_{\text{甲}}}{\dfrac{1}{3}\left( S_2+S_1+\sqrt{S_2S_1} \right)h_{\text{乙}}}
=\frac{h_{\text{甲}}}{h_{\text{乙}}}
=\frac{\sqrt{3}\left( r_1-r_2 \right)}{2\sqrt{2}\left( r_1-r_2 \right)}
=\frac{\sqrt{6}}{4}
\]

答案为：\( \dfrac{\sqrt{6}}{4} \).
\end{solution}

\begin{problem}
已知 \( a>1 \)，\( \dfrac{1}{\log_8 a}-\dfrac{1}{\log_a 4}=-\dfrac{5}{2} \)，则 \( a= \)\fillinblank{}.

\end{problem}

\begin{answer}
64
\end{answer}

\begin{solution}
由题
\[
\frac{1}{\log_8 a}-\frac{1}{\log_a 4}=\frac{3}{\log_2 a}-\frac{1}{2}\log_2 a=-\frac{5}{2}
\]
整理得
\[
\left( \log_2 a \right)^2-5\log_2 a-6=0
\]

所以
\(\log_2 a=-1\)
或
\(\log_2 a=6\).

又 \( a>1 \)，所以
\(\log_2 a=6=\log_2 2^6\)，
故
\(a=2^6=64\).

答案为：64.
\end{solution}

\begin{problem}
有 \( 6 \) 个相同的球，分别标有数字 \( 1,2,3,4,5,6 \)，从中不放回地随机抽取 \( 3 \) 次，每次取 \( 1 \) 个球. 记 \( m \) 为前两次取出的球上数字的平均值，\( n \) 为取出的三个球上数字的平均值，则 \( m \) 与 \( n \) 差的绝对值不超过 \( \dfrac{1}{2} \) 的概率是\fillinblank{}.

\end{problem}

\begin{answer}
\( \dfrac{7}{15} \)
\end{answer}

\begin{solution}
从 \( 6 \) 个不同的球中不放回地抽取 \( 3 \) 次，共有
\(A_6^3=120\)
种.

设前两个球的号码为 \( a,b \)，第三个球的号码为 \( c \)，则
\(\left| \frac{a+b+c}{3}-\frac{a+b}{2} \right|\le \frac{1}{2}\)，
所以
\(\left| 2c-(a+b) \right|\le 3\)，
即
\(a+b-3\le 2c\le a+b+3\).

若 \( c=1 \)，则 \( a+b\le 5 \)，则 \( (a,b) \) 为 \((2,3)\)，\((3,2)\)，因此有 \( 2 \) 种.

若 \( c=2 \)，则 \( 1\le a+b\le 7 \)，则 \( (a,b) \) 为
\((1,3)\)，\((1,4)\)，\((1,5)\)，\((1,6)\)，\((3,4)\)，
\((3,1)\)，\((4,1)\)，\((5,1)\)，\((6,1)\)，\((4,3)\)，
因此有 \( 10 \) 种.

当 \( c=3 \) 时，则 \( 3\le a+b\le 9 \)，则 \( (a,b) \) 为
\((1,2)\)，\((1,4)\)，\((1,5)\)，\((1,6)\)，\((2,4)\)，\((2,5)\)，\((2,6)\)，\((4,5)\)，
\((2,1)\)，\((4,1)\)，\((5,1)\)，\((6,1)\)，\((4,2)\)，\((5,2)\)，\((6,2)\)，\((5,4)\)，
因此有 \( 16 \) 种.

当 \( c=4 \) 时，则 \( 5\le a+b\le 11 \)，则 \( (a,b) \) 为
\((1,5)\)，\((1,6)\)，\((2,3)\)，\((2,5)\)，\((2,6)\)，\((3,2)\)，\((3,5)\)，\((3,6)\)，
\((5,1)\)，\((5,2)\)，\((5,3)\)，\((5,6)\)，\((6,1)\)，\((6,2)\)，\((6,3)\)，\((6,5)\)，
因此有 \( 16 \) 种.

当 \( c=5 \) 时，则 \( 7\le a+b\le 13 \)，则 \( (a,b) \) 为
\((1,6)\)，\((2,6)\)，\((3,4)\)，\((3,6)\)，\((4,3)\)，
\((4,6)\)，\((6,1)\)，\((6,2)\)，\((6,3)\)，\((6,4)\)，
因此有 \( 10 \) 种.

当 \( c=6 \) 时，则 \( 9\le a+b\le 15 \)，则 \( (a,b) \) 为 \((4,5)\)，\((5,4)\)，因此有 \( 2 \) 种.

共 \( m \) 与 \( n \) 的差的绝对值不超过 \( \dfrac{1}{2} \) 时不同的抽取方法总数为
\(2(2+10+16)=56\)，
因此所求概率为
\(\frac{56}{120}=\frac{7}{15}\).

答案为：\( \dfrac{7}{15} \).
\end{solution}
\section{解答题}

\begin{problem}
某工厂进行生产线智能化升级改造，升级改造后，从该工厂甲、乙两个车间的产品中随机抽取 \( 150 \) 件进行检验，数据如下：

\begin{center}
\renewcommand{\arraystretch}{1.3}
\begin{tabular}{c|c|c|c|c}
 & 优级品 & 合格品 & 不合格品 & 总计 \\
\hline
甲车间 & 26 & 24 & 0 & 50 \\
\hline
乙车间 & 70 & 28 & 2 & 100 \\
\hline
总计 & 96 & 52 & 2 & 150
\end{tabular}
\end{center}

\begin{enumerate}
\item 填写如下列联表：

\begin{center}
\renewcommand{\arraystretch}{1.3}
\begin{tabular}{c|c|c}
 & 优级品 & 非优级品 \\
\hline
甲车间 &  &  \\
\hline
乙车间 &  &
\end{tabular}
\end{center}

能否有 \( 95\% \) 的把握认为甲、乙两车间产品的优级品率存在差异？能否有 \( 99\% \) 的把握认为甲、乙两车间产品的优级品率存在差异？
\item 已知升级改造前该工厂产品的优级品率 \( p=0.5 \)，设 \( \bar p \) 为升级改造后抽取的 \( n \) 件产品的优级品率. 如果 \( \bar p>p+1.65\sqrt{\frac{p(1-p)}{n}} \)，则认为该工厂产品的优级品率提高了. 根据抽取的 \( 150 \) 件产品的数据，能否认为生产线智能化升级改造后，该工厂产品的优级品率提高了？（\( \sqrt{150}\approx 12.247 \)）
\end{enumerate}

附：\( K^2=\frac{n(ad-bc)^2}{(a+b)(c+d)(a+c)(b+d)} \)

\begin{center}
\renewcommand{\arraystretch}{1.3}
\begin{tabular}{c|c|c|c}
\( P\left( K^2\ge k \right) \) & 0.050 & 0.010 & 0.001 \\
\hline
\( k \) & 3.841 & 6.635 & 10.828
\end{tabular}
\end{center}

\end{problem}

\begin{solution}
（1）根据题意可得列联表：

\begin{center}
\renewcommand{\arraystretch}{1.3}
\begin{tabular}{c|c|c}
 & 优级品 & 非优级品 \\
\hline
甲车间 & 26 & 24 \\
\hline
乙车间 & 70 & 30
\end{tabular}
\end{center}

可得
\[
K^2=\frac{150\left( 26\times 30-24\times 70 \right)^2}{50\times 100\times 96\times 54}=\frac{75}{16}=4.6875
\]
因为
\(3.841<4.6875<6.635\)，
所以有 \( 95\% \) 的把握认为甲、乙两车间产品的优级品率存在差异，没有 \( 99\% \) 的把握认为甲、乙两车间产品的优级品率存在差异.

（2）根据题意可知：生产线智能化升级改造后，该工厂产品的优级品的频率为
\(\frac{96}{150}=0.64\).

用频率估计概率可得 \( \bar p=0.64 \).

又因为升级改造前该工厂产品的优级品率 \( p=0.5 \)，则
\[
p+1.65\sqrt{\frac{p(1-p)}{n}}
=0.5+1.65\sqrt{\frac{0.5(1-0.5)}{150}}
\approx 0.5+1.65\times \frac{0.5}{12.247}
\approx 0.568
\]
可知
\(\bar p>p+1.65\sqrt{\frac{p(1-p)}{n}}\)，
所以可以认为生产线智能化升级改造后，该工厂产品的优级品率提高了.
\end{solution}

\begin{problem}
记 \( S_n \) 为数列 \( \{a_n\} \) 的前 \( n \) 项和，且 \( 4S_n=3a_n+4 \).

\begin{enumerate}
\item 求 \( \{a_n\} \) 的通项公式；
\item 设 \( b_n=\left( -1 \right)^{n-1}na_n \)，求数列 \( \{b_n\} \) 的前 \( n \) 项和 \( T_n \).
\end{enumerate}

\end{problem}

\begin{solution}
（1）当 \( n=1 \) 时，
\[
4S_1=4a_1=3a_1+4
\]
解得 \( a_1=4 \).

当 \( n\ge 2 \) 时，
\[
4S_{n-1}=3a_{n-1}+4
\]
所以
\[
4S_n-4S_{n-1}=4a_n=3a_n-3a_{n-1}
\]
即 \( a_n=-3a_{n-1} \).

而 \( a_1=4\ne 0 \)，故 \( a_n\ne 0 \)，所以
\(\frac{a_n}{a_{n-1}}=-3\).

故数列 \( \{a_n\} \) 是以 \( 4 \) 为首项，\( -3 \) 为公比的等比数列，
\(a_n=4\cdot \left( -3 \right)^{n-1}\).

（2）
\(b_n=\left( -1 \right)^{n-1}n\cdot 4\cdot \left( -3 \right)^{n-1}=4n\cdot 3^{n-1}\)，
所以
\(T_n=b_1+b_2+b_3+\cdots+b_n =4\cdot 3^0+8\cdot 3^1+12\cdot 3^2+\cdots+4n\cdot 3^{n-1}\).

则
\(3T_n=4\cdot 3^1+8\cdot 3^2+12\cdot 3^3+\cdots+4n\cdot 3^n\)，
所以
\[
\begin{aligned}
-2T_n
&=4+4\cdot 3^1+4\cdot 3^2+\cdots+4\cdot 3^{n-1}-4n\cdot 3^n \\
&=4+4\cdot \frac{3\left( 1-3^{n-1} \right)}{1-3}-4n\cdot 3^n \\
&=4+2\cdot 3\left( 3^{n-1}-1 \right)-4n\cdot 3^n \\
&=\left( 2-4n \right)\cdot 3^n-2.
\end{aligned}
\]

故
\(T_n=\left( 2n-1 \right)\cdot 3^n+1\).
\end{solution}

\begin{problem}
如图，在以 \( A,B,C,D,E,F \) 为顶点的五面体中，四边形 \( ABCD \) 与四边形 \( ADEF \) 均为等腰梯形，\( BC//AD \)，\( EF//AD \)，\( AD=4 \)，\( AB=BC=EF=2 \)，\( ED=\sqrt{10} \)，\( FB=2\sqrt{3} \)，\( M \) 为 \( AD \) 的中点.

\begin{enumerate}
\item 证明：\( BM// \)平面 \( CDE \)；
\item 求二面角 \( F-BM-E \) 的正弦值.
\end{enumerate}

\begin{center}
\bitmapinclude[max width=0.55\linewidth,max totalheight=4.8cm]{img/2024/national_paper_A_science/q19_fig1.png}
\end{center}
\end{problem}

\begin{solution}
（1）因为 \( BC//AD \)，\( BC=2 \)，\( AD=4 \)，\( M \) 为 \( AD \) 的中点，所以 \( BC//MD \)，\( BC=MD \).

四边形 \( BCDM \) 为平行四边形，所以 \( BM//CD \)，又因为 \( BM\nsubseteq \) 平面 \( CDE \)，\( CD\subset \) 平面 \( CDE \)，所以 \( BM// \)平面 \( CDE \).

（2）作 \( BO\perp AD \) 交 \( AD \) 于 \( O \)，连接 \( OF \).

因为四边形 \( ABCD \) 为等腰梯形，\( BC//AD \)，\( AD=4 \)，\( AB=BC=2 \)，所以 \( CD=2 \).

结合（1）\( BCDM \) 为平行四边形，可得 \( BM=CD=2 \)，又 \( AM=2 \)，所以 \( \triangle ABM \) 为等边三角形，\( O \) 为 \( AM \) 中点，所以
\(OB=\sqrt{3}\).

又因为四边形 \( ADEF \) 为等腰梯形，\( M \) 为 \( AD \) 中点，所以 \( EF=MD \)，\( EF//MD \).

四边形 \( EFMD \) 为平行四边形，\( FM=ED=AF \)，所以 \( \triangle AFM \) 为等腰三角形，\( \triangle ABM \) 与 \( \triangle AFM \) 底边上中点 \( O \) 重合，且 \( OF\perp AM \)，
\(OF=\sqrt{AF^2-AO^2}=3\).

因为
\(OB^2+OF^2=BF^2\)，
所以 \( OB\perp OF \)，故 \( OB,OD,OF \) 互相垂直.

以 \( OB \) 方向为 \( x \) 轴，\( OD \) 方向为 \( y \) 轴，\( OF \) 方向为 \( z \) 轴，建立 \( O-xyz \) 空间直角坐标系.

则
\(F(0,0,3)\)，\(B(3,0,0)\)，\(M(0,1,0)\)，\(E(0,2,3)\)，
所以
\(\overrightarrow{BM}=(-3,1,0)\)，\(\overrightarrow{BF}=(-3,0,3)\)，\(\overrightarrow{BE}=(-3,2,3)\).

设平面 \( BFM \) 的法向量为 \( \bs m=(x_1,y_1,z_1) \)，平面 \( EMB \) 的法向量为 \( \bs n=(x_2,y_2,z_2) \).

则 \( \bs m\cdot \overrightarrow{BM}=0 \)，\( \bs m\cdot \overrightarrow{BF}=0 \)，即 \( -3x_1+y_1=0 \)，\( -3x_1+3z_1=0 \). 令 \( x_1=3 \)，得 \( y_1=3 \)，\( z_1=1 \)，即
\(\bs m=(3,3,1)\).

又因为 \( \bs n \) 为平面 \( EMB \) 的法向量，所以 \( \bs n\cdot \overrightarrow{BM}=0 \)，\( \bs n\cdot \overrightarrow{BE}=0 \)，即 \( -3x_2+y_2=0 \)，\( -3x_2+2y_2+3z_2=0 \). 令 \( x_2=3 \)，得 \( y_2=3 \)，\( z_2=-1 \)，即
\(\bs n=(3,3,-1)\).

于是
\(\cos \angle \left( \bs m,\bs n \right)=\frac{\bs m\cdot \bs n}{|\bs m||\bs n|}=\frac{11}{13}\)，
所以
\(\sin \angle \left( \bs m,\bs n \right)=\frac{4\sqrt{3}}{13}\).

故二面角 \( F-BM-E \) 的正弦值为 \( \dfrac{4\sqrt{3}}{13} \).
\end{solution}

\begin{problem}
设椭圆 \( C:\dfrac{x^2}{a^2}+\dfrac{y^2}{b^2}=1 \)（\( a>b>0 \)）的右焦点为 \( F \)，点 \( M\left( 1,\dfrac{\sqrt{3}}{2} \right) \) 在 \( C \) 上，且 \( MF\perp x \) 轴.

\begin{enumerate}
\item 求 \( C \) 的方程；
\item 过点 \( P(4,0) \) 的直线与 \( C \) 交于 \( A,B \) 两点，\( N \) 为线段 \( FP \) 的中点，直线 \( NB \) 交直线 \( MF \) 于点 \( Q \)，证明：\( AQ\perp y \) 轴.
\end{enumerate}

\begin{center}
\bitmapinclude[max width=0.62\linewidth,max totalheight=4.8cm]{img/2024/national_paper_A_science/q20_fig2.png}
\end{center}
\end{problem}

\begin{solution}
（1）设 \( F(c,0) \)，由题设有 \( c=1 \) 且 \( \frac{b^2}{a^2}=\frac{3}{4} \)，故
\[
\frac{a^2-1}{a^2}=\frac{3}{4}
\]
于是 \( a=2 \)，故 \( b=\sqrt{3} \).

所以椭圆方程为
\(\frac{x^2}{4}+\frac{y^2}{3}=1\).

（2）直线 \( AB \) 的斜率必定存在，设
\(AB:y=k(x-4)\)，
其中 \( A(x_1,y_1) \)，\( B(x_2,y_2) \).

由 \( 3x^2+4y^2=12 \) 与 \( y=k(x-4) \) 可得
\[
\left( 3+4k^2 \right)x^2-32k^2x+64k^2-12=0
\]
故
\[
\Delta=1024k^4-4\left( 3+4k^2 \right)\left( 64k^2-12 \right)>0
\]
从而 \( -\frac{1}{2}<k<\frac{1}{2} \).

又
\(x_1+x_2=\frac{32k^2}{3+4k^2}\)，\(x_1x_2=\frac{64k^2-12}{3+4k^2}\).

因为 \( N\left( \dfrac{5}{2},0 \right) \)，故直线 \( BN \) 的方程为
\(y=\frac{y_2}{x_2-\frac{5}{2}}\left( x-\frac{5}{2} \right)\).

由 \( MF\perp x \) 轴且 \( F(1,0) \)，可知直线 \( MF \) 的方程为 \( x=1 \)，故
\(y_Q=\frac{-3y_2}{2x_2-5}\).

所以
\[
\begin{aligned}
y_1-y_Q
&=y_1+\frac{3y_2}{2x_2-5}\\
&=\frac{y_1(2x_2-5)+3y_2}{2x_2-5}\\
&=\frac{k(x_1-4)(2x_2-5)+3k(x_2-4)}{2x_2-5}\\
&=\frac{k\left( 2x_1x_2-5x_1+x_2-8+3x_2-12 \right)}{2x_2-5}\\
&=\frac{k\left( 2x_1x_2-5(x_1+x_2)+8 \right)}{2x_2-5}.
\end{aligned}
\]

代入 \( x_1+x_2 \) 与 \( x_1x_2 \) 得
\[
\begin{aligned}
2x_1x_2-5(x_1+x_2)+8
&=2\times \frac{64k^2-12}{3+4k^2}-5\times \frac{32k^2}{3+4k^2}+8\\
&=\frac{128k^2-24-160k^2+24+32k^2}{3+4k^2}\\
&=0.
\end{aligned}
\]

故
\(y_1=y_Q\)，
即 \( AQ\perp y \) 轴.
\end{solution}

\begin{problem}
已知函数 \( f(x)=(1-ax)\ln(1+x)-x \).

\begin{enumerate}
\item 当 \( a=-2 \) 时，求 \( f(x) \) 的极值；
\item 当 \( x\ge 0 \) 时，\( f(x)\ge 0 \) 恒成立，求 \( a \) 的取值范围.
\end{enumerate}

\end{problem}

\begin{solution}
（1）当 \( a=-2 \) 时，\( f(x)=(1+2x)\ln(1+x)-x \).

故
\[
f'(x)=2\ln(1+x)+\frac{1+2x}{1+x}-1=2\ln(1+x)-\frac{1}{1+x}+1
\]

因为
\(y=2\ln(1+x)\)，\(y=-\frac{1}{1+x}+1\)
在 \( (-1,+\infty) \) 上都为增函数，故 \( f'(x) \) 在 \( (-1,+\infty) \) 上为增函数，而
\(f'(0)=0\).

故当 \( -1<x<0 \) 时，\( f'(x)<0 \)；当 \( x>0 \) 时，\( f'(x)>0 \).

所以 \( f(x) \) 在 \( x=0 \) 处取极小值，且极小值为
\(f(0)=0\)，
无极大值.

（2）
\(f'(x)=-a\ln(1+x)+\frac{1-ax}{1+x}-1=-a\ln(1+x)-\frac{(a+1)x}{1+x}\)，\(x>0\).

设
\(s(x)=-a\ln(1+x)-\frac{(a+1)x}{1+x}\)，\(x>0\)，
则
\[
s'(x)=-\frac{a}{1+x}-\frac{a+1}{(1+x)^2}
=-\frac{a(x+1)+a+1}{(1+x)^2}
=-\frac{ax+2a+1}{(1+x)^2}
\]

当 \( a\le -\dfrac{1}{2} \) 时，\( s'(x)>0 \)，故 \( s(x) \) 在 \( (0,+\infty) \) 上为增函数，
\(s(x)>s(0)=0\)，
即 \( f'(x)>0 \)，所以 \( f(x) \) 在 \( [0,+\infty) \) 上为增函数，故
\(f(x)\ge f(0)=0\).

当 \( -\dfrac{1}{2}<a<0 \) 时，当
\(0<x<-\frac{2a+1}{a}\)
时，\( s'(x)<0 \)，故 \( s(x) \) 在 \( \left( 0,-\dfrac{2a+1}{a} \right) \) 上为减函数，
\(s(x)<s(0)\)，
即在 \( \left( 0,-\dfrac{2a+1}{a} \right) \) 上 \( f'(x)<0 \)，所以 \( f(x)<f(0)=0 \)，不合题意.

当 \( a\ge 0 \) 时，此时 \( s'(x)<0 \) 在 \( (0,+\infty) \) 上恒成立，故 \( s(x) \) 在 \( (0,+\infty) \) 上为减函数，
\(s(x)<s(0)=0\)，
即 \( f'(x)<0 \)，所以 \( f(x) \) 在 \( (0,+\infty) \) 上为减函数，从而
\(f(x)<f(0)=0\)，
也不合题意.

综上，
\(a\le -\frac{1}{2}\).
\end{solution}

\begin{problem}
在平面直角坐标系 \( xOy \) 中，以坐标原点 \( O \) 为极点，\( x \) 轴的正半轴为极轴建立极坐标系，曲线 \( C \) 的极坐标方程为 \( \rho=\rho\cos\theta+1 \).

\begin{enumerate}
\item 写出 \( C \) 的直角坐标方程；
\item 设直线 \( l \) 的参数方程为 \( x=t \)，\( y=t+a \)（\( t \) 为参数），若 \( C \) 与 \( l \) 相交于 \( A,B \) 两点，若 \( |AB|=2 \)，求 \( a \) 的值.
\end{enumerate}

\end{problem}

\begin{solution}
（1）由 \( \rho=\rho\cos\theta+1 \)，将 \( \rho=\sqrt{x^2+y^2} \) 和 \( \rho\cos\theta=x \) 代入，可得
\(\sqrt{x^2+y^2}=x+1\).

两边平方后可得曲线的直角坐标方程为
\(y^2=2x+1\).

（2）对于直线 \( l \) 的参数方程消去参数 \( t \)，得直线的普通方程为
\(y=x+a\).

法 1：直线 \( l \) 的斜率为 \( 1 \)，故倾斜角为 \( \dfrac{\pi}{4} \)，故直线的参数方程可设为 \( x=\dfrac{\sqrt{2}}{2}s \)，\( y=a+\dfrac{\sqrt{2}}{2}s \)（\( s\in \R \)）.

将其代入 \( y^2=2x+1 \) 中得
\(s^2+2\sqrt{2}(a-1)s+2(a^2-1)=0\).

设 \( A,B \) 两点对应的参数分别为 \( s_1,s_2 \)，则
\(s_1+s_2=-2\sqrt{2}(a-1)\)，\(s_1s_2=2(a^2-1)\).

且
\[
\Delta=8(a-1)^2-8(a^2-1)=16-16a>0
\]
故 \( a<1 \).

又
\[
|AB|=|s_1-s_2|=\sqrt{(s_1+s_2)^2-4s_1s_2}=\sqrt{8(a-1)^2-8(a^2-1)}=2
\]
解得 \( a=\frac{3}{4} \).

法 2：联立 \( y=x+a \) 与 \( y^2=2x+1 \)，得
\(x^2+(2a-2)x+a^2-1=0\).

则
\[
\Delta=(2a-2)^2-4(a^2-1)=-8a+8>0
\]
解得 \( a<1 \). 设 \( A(x_1,y_1) \)，\( B(x_2,y_2) \)，则
\(x_1+x_2=2-2a\)，\(x_1x_2=a^2-1\).

所以
\[
|AB|=\sqrt{1+1^2}\cdot \sqrt{(x_1+x_2)^2-4x_1x_2}=\sqrt{2}\cdot \sqrt{(2-2a)^2-4(a^2-1)}=2
\]
解得 \( a=\frac{3}{4} \).
\end{solution}

\begin{problem}
实数 \( a,b \) 满足 \( a+b\ge 3 \).

\begin{enumerate}
\item 证明：\( 2a^2+2b^2>a+b \)；
\item 证明：\( |a-2b^2|+|b-2a^2|\ge 6 \).
\end{enumerate}

\end{problem}

\begin{solution}
（1）因为
\[
2a^2+2b^2-(a+b)^2=a^2-2ab+b^2=(a-b)^2\ge 0
\]
当 \( a=b \) 时等号成立，则
\(2a^2+2b^2\ge (a+b)^2\).

因为 \( a+b\ge 3 \)，所以
\(2a^2+2b^2\ge (a+b)^2>a+b\).

（2）
\[
\begin{aligned}
|a-2b^2|+|b-2a^2|
&\ge |a-2b^2+b-2a^2|\\
&=|2a^2+2b^2-(a+b)|\\
&=2a^2+2b^2-(a+b)\\
&\ge (a+b)^2-(a+b)\\
&=(a+b)(a+b-1)\ge 3\times 2=6.
\end{aligned}
\]

故
\(|a-2b^2|+|b-2a^2|\ge 6\).
\end{solution}
